\clearpage
\renewcommand{\Article}[1]{\par\medskip\Needspace{4\baselineskip}\pdfbookmark[1]{Art. #1}{nexus-second-prior-art-#1}\centerline{\textbf{#1.}}\nobreak\smallskip}
\pdfbookmark[0]{The Number of Classes and Their Determinant. II}{nexus-second-prior}
\begin{center}
[II.]\par\medskip
\Large\bfseries ON THE CONNECTION BETWEEN THE NUMBER OF CLASSES INTO WHICH BINARY QUADRATIC FORMS ARE DISTRIBUTED AND THEIR DETERMINANT.\par\medskip
\normalsize FIRST PAPER\par\medskip
PRESENTED TO THE ROYAL SOCIETY 1837\ .\ .\ .
\end{center}
\BookRule{.18\textwidth}

\Article{1}
Thirty-six years have passed since the principles of the remarkable connection to be treated in this paper were discovered, as was already announced at the end of the \emph{Arithmetical Investigations}. Other occupations, however, had long drawn attention away from this inquiry, until more recently it became possible to return to it and extend it through renewed efforts. Yet, since this new part of higher arithmetic exceeds the bounds of a single paper, this first paper will be devoted to forms with negative determinants; forms with positive determinants, which require an entirely distinct treatment, will remain reserved for another paper.

\SourceLeaf{287}
\Article{2}
For our purpose we shall need a theorem which is itself arithmetical, but whose nature can be presented more conveniently and clearly through considerations expressed in geometrical form.

Given a figure in an unbounded plane, enclosed by any line, its area can be determined approximately by dividing the plane into squares and counting both those lying wholly inside the figure and those cut by its boundary. Evidently, the resulting area will be smaller or larger than the true area according as the latter squares are omitted or added to the former. If, however, we choose, by any rule, to exclude some and include others of the latter squares lying on the boundary, the error may be positive or negative, but must be smaller than the total area of all the boundary squares. The smaller the squares are taken, the more accurately the area will be determined in this way. We may carry such an approximation indefinitely further, or take the squares so small that the error becomes smaller than any given quantity. Although this may already seem evident in itself, we shall nevertheless not refuse to support it with a rigorous proof.

\Needspace{6\baselineskip}
Two squares can have one common vertex, two, or none; in the first and second cases they will be called adjacent, and in the third disjoint. Evidently, at most four squares can all be adjacent to one another, and thus among any five distinct squares at least two must be disjoint. Now the distance between two points lying in disjoint squares cannot be smaller than the side of the squares, which we shall denote by $a$. Hence, if a point starts anywhere in one square, then passes through a second, third and fourth, and finally reaches a fifth, the length of its path certainly cannot be smaller than $a$. By the same reasoning, if the line continues through further squares, the part between the fifth and ninth, and likewise between the ninth and thirteenth, etc., cannot be smaller than $a$. We therefore readily conclude that any line returning to itself which has touched $n$ distinct squares in all certainly cannot be shorter than $\frac{(n-4)a}4$. Conversely, a closed line whose length is $=l$ certainly cannot have touched more than $4+\frac{4l}a$ distinct squares. Since their area $=4aa+4al$ can become smaller than any given quantity as $a$ decreases indefinitely, the same holds all the more for the error in the determination of area discussed above.

\SourceLeaf{288}
\Article{3}
The rule for including or excluding squares lying on the boundary of the figure could be chosen in many different ways. The simplest, however, seems to be to consider only the position of the centre of each square, so that squares whose centres lie inside the figure are included, those whose centres lie outside it are excluded, and it is left to our choice whether centres that happen to lie on the boundary itself are counted with the interior or the exterior. Instead of the centres, we could also take any other points situated in the same relative position in each square.

Thus the matter comes down to imagining equally spaced points distributed along equally spaced straight lines in the plane so as to form squares. By the theorem of the preceding article, we can then assert that the number of points contained in the figure, multiplied by the square of the distance between two nearest points, becomes as nearly equal to the area of the figure as desired, provided that distance is taken sufficiently small; or, in the customary manner of speaking, that this product gives the area if the distance is infinitely small.

\Article{4}
The curve expressed by the following equation between the rectangular coordinates $p,q$,
\[
app+2bpq+cqq=1
\]
is a conic section, and specifically an ellipse if $a,c$ and $ac-bb$ are positive quantities. The area enclosed by this ellipse is found to be $\pdot=\frac\pi{\surd(ac-bb)}$. \mbox{The value of $app+2bpq+cqq$} is everywhere greater than $1$ outside the ellipse, less than $1$ inside it, and nowhere negative.

Suppose a .system of points distributed in the plane containing the ellipse so as to form squares whose sides $=\lambda$ are parallel to the coordinate axes. It makes no difference whether the origin of the coordinates, or centre of the ellipse, coincides with one of these points. Let the number of points inside the ellipse, including any on its boundary, be $=m$. By the theorem of the preceding article, $\frac\pi{\surd(ac-bb)}$ will then be the limit of $m\lambda\lambda$, which approaches it as closely as desired as $\lambda$ decreases indefinitely.

If we suppose that the origin of the coordinates coincides with a point of the system, then, on putting $p=\lambda x,q=\lambda y$, evidently $x$ and $y$ will be integers for each point of the system; conversely, every combination of integer values of
\SourceLeaf{289}
$x,y$ will correspond to a point of the system. Thus the number $m$ is simply the number of all combinations of integer values of $x,y$ for which $F$ is no greater than $M$, if, for brevity, we denote the function, or quadratic form, $axx+2bxy+cyy$ by $F$, and the quantity $\frac1{\lambda\lambda}$ by $M$. The determinant of this form is $bb-ac$, for which we shall write $-D$. Our theorem can now be stated as follows.

\textsc{Theorem I.}\quad\emph{The number $m$ of all combinations of integer values of the indeterminates $x,y$ for which the value of a form with negative determinant $-D$ does not exceed the bound $M$ becomes $=\frac{\pi M}{\surd D}$, approximately, indeed, but with the approximation improving indefinitely as $M$ increases indefinitely. It scarcely needs to be mentioned that indefinite approximation here, and likewise in what follows, is not to be understood as saying that the difference between $\frac{\pi M}{\surd D}$ and $m$ itself decreases indefinitely. Rather, the ratio between these quantities will approach equality indefinitely; that is, $\frac{\pi M}{m\surd D}-1$ will decrease indefinitely.}

\Article{5}
To carry out the count, we may proceed by computing, for each integer value of $x$ between the bounds $-\surd\frac{cM}D$ and $+\surd\frac{cM}D$, the two values of $y$ satisfying $F=M$; the number of integers lying between them is then obtained directly. Since this number is the same for opposite values of $x$, almost half the work is saved. The count can also be carried out by counting the values of $x$ corresponding to each value of $y$ between the bounds $-\surd\frac{aM}D$ and $+\surd\frac{aM}D$. A suitable combination of the two methods can reduce the work further, but we do not pursue this here at greater length; it will suffice to add a few observations on the simplest case.

\Needspace{12\baselineskip}
Let the form be $F=xx+yy$, so that the curve is a circle, and let $r,r',r'',r'''\ldots r^{(r)}$ denote the integers immediately below
\[
\surd M,\quad\surd(M-1),\quad\surd(M-4),\quad\surd(M-9)\ldots\surd(M-rr)
\]
or these quantities themselves, if any of them are integers. The required number will then be
\[
\begin{aligned}
m&=2r+1+2(2r'+1)+2(2r''+1)+2(2r'''+1)+\text{etc.}+2(2r^{(r)}+1)\\
&=1+4r+4r'+4r''+4r'''+\text{etc.}+4r^{(r)}
\end{aligned}
\]

We obtain the same result more quickly, however, by denoting by $q$ the integer immediately
\SourceLeaf{290}
below $\surd\tfrac12 M$ (or this quantity itself, if it is an integer), with the aid of the formula
\[
m=4qq+1+4r+8(r^{(q+1)}+r^{(q+2)}+r^{(q+3)}+\text{etc.}+r^{(r)})
\]
The following numbers were obtained in this way.
\[
\begin{array}{r|r||r|r||r|r}
M&m&M&m&M&m\\\hline
100&317&1000&3149&10000&31417\\
200&633&2000&6293&20000&62845\\
300&949&3000&9425&30000&94237\\
400&1257&4000&12581&40000&125629\\
500&1581&5000&15705&50000&157093\\
600&1885&6000&18853&60000&188453\\
700&2209&7000&21993&70000&219901\\
800&2521&8000&25137&80000&251305\\
900&2821&9000&28269&90000&282697\\
1000&3149&10000&31417&100000&314197
\end{array}
\]

\Article{6}
We give the theorem of Article 4 greater generality in the following way.

\textsc{Theorem II.}\quad\emph{If we are to collect, not all combinations of integer values of $x,y$ for which $F$ does not exceed the value $M$, but only those taken at intervals, namely those for which $x$ is congruent to a given number $G$ to a given modulus $g$, and $y$ is congruent to a given number $H$ to a given modulus $h$, then the number $m'$ of these combinations will be expressed approximately by $\frac{\pi M}{gh\surd D}$, with the approximation improving indefinitely as $M$ increases indefinitely.}

\Needspace{7\baselineskip}
Indeed, on putting $x=gx'+G,y=hy'+H$, it is evident that $m'$ is the number of all combinations of integer values of $x',y'$ for which
\begingroup\postdisplaypenalty=10000
\[
agg(x'+\tfrac Gg)^2+2bgh(x'+\tfrac Gg)(y'+\tfrac Hh)+chh(y'+\tfrac Hh)^2
\]
does not exceed the value $M$.
\endgroup
Thus, if we suppose a system of points distributed in the plane as in Article 4, but so that the point with coordinates $p=\frac{G\lambda}g,q=\frac{H\lambda}h$, rather than the origin of the coordinates, coincides with a point of the system, then $m'$ will express the number of points lying inside the
\SourceLeaf{291}
ellipse whose equation is
\begingroup\postdisplaypenalty=10000
\[
aggpp+2bghpq+chhqq=1
\]
with any points on the boundary itself always included.
\endgroup The area of this ellipse $=\frac\pi{gh\surd(ac-bb)}=\frac\pi{gh\surd D}$ will be the limit which the product $m'\lambda\lambda=\frac{m'}M$ approaches indefinitely as $\lambda$ decreases or $M$ increases indefinitely.

It is also evident that our theorem includes the case where only one of the indeterminates $x,y$ must proceed by intervals, while the value of the other is subject to no condition. Clearly, this is the same as putting either $h$ or $g$ $=1$.

\Article{7}
What has been set out so far is independent of the nature of the coefficients of the form $axx+2bxy+cyy$; from now on, however, we shall suppose these coefficients to be integers. Thus every combination of integer values of $x,y$ gives the form itself an integer value, or corresponds to a representation of some integer by this form. Hence the collection of all combinations of integer values of $x,y$ for which the form $F=axx+2bxy+cyy$ takes a value no greater than the bound $M$ is the same as the collection of all representations of integers not exceeding the bound $M$, or up to and including this bound if it is itself an integer. If, for brevity, we denote the number of distinct representations of a specified integer $n$ by the form $F$ by $F(n)$, or simply by $Fn$ when there is no risk of ambiguity, then the number previously expressed by $m$ will be $=F0+F1+F2+F3+\text{etc.}+FM$, and the first theorem takes the following form.

\textsc{Theorem III.}\quad\emph{The sum $F0+F1+F2+\text{etc.}+FM$ is expressed approximately by $\frac{\pi M}{\surd D}$, with the approximation improving indefinitely as $M$ increases indefinitely.}

\Article{8}
To the third theorem, which concerns representations of \emph{all} numbers, it is useful to add another concerning odd numbers alone. Evidently, odd numbers cannot be represented by the form $F$ if $a$ and $c$ are both even numbers; the investigation will therefore be restricted to the three remaining cases.

I. Whenever $a$ is odd and $c$ is even, an odd number is represented by giving $x$ an odd value, while the value of $y$ remains arbitrary. Theorem II,
\SourceLeaf{292}
on putting $g=2,G=1,h=1$, therefore shows that the number of all combinations of such values of $x,y$ as give the form an odd value no greater than the bound $M$ is expressed, with indefinite approximation, by $\frac{\pi M}{2\surd D}$, as $M$ increases indefinitely.

II. Whenever $a$ is even and $c$ is odd, the representation of an odd number requires $y$ to be odd. Thus, on putting $g=1,h=2,H=1$, we reach the same conclusion.

\Needspace{8\baselineskip}
III. Whenever both $a$ and $c$ are odd, either an odd value of $x$ must be combined with an even value of $y$, or an even value of $x$ with an odd value of $y$, for the formula to produce an odd value. The number of all combinations of each kind for which the value of the form does not exceed the bound $M$ is expressed, with indefinite approximation, by $\frac{\pi M}{4\surd D}$. Therefore the number of all combinations which give the form odd values not exceeding the bound $M$ is again expressed, with indefinite approximation, by $\frac{\pi M}{2\surd D}$.

Now the collection of all such combinations is simply the collection of all representations of all the numbers $1,3,5,7\ldots M$ when $M$ is an odd integer, or $1,3,5,7\dotsfour M-1$ when $M$ is even. We therefore have

\textsc{Theorem IV.}\quad\emph{The sum}
\[
F1+F3+F5+F7\ldotp\ldotp+FM\quad\text{or}\quad F1+F3+F5+F7\ldotp\ldotp+F(M-1)
\]
\emph{(according as $M$ is odd or even) is expressed, with indefinite approximation, by $\frac{\pi M}{2\surd D}$, provided that $F$ is a form in which one or both of the coefficients $a,c$ are odd.}
\BookRule{.30\textwidth}
